% §6.3 Cancellation and the Loss of One-Sidedness
%
% Monograph register + "B" fix continuation. The mixed-axis intersection
% (FNR-zero cap FPR-zero) is two-sided IMMEDIATELY by the corrected
% eq:intersection-fpr (the w2*eps_A term survives), not landing at (0, omega_B);
% prop:mixed-intersection-escape restated accordingly. Cross-refs def:one-sided,
% sec:union-of-sets, prop:union-fpr-fnr, sec:other-set-ops,
% prop:intersection-fpr-fnr, eq:intersection-fpr, eq:intersection-fnr,
% prop:complement-swap, prop:difference-fpr-fnr, sec:perfect-filter-line,
% ch:classification-measures, ch:optimal-bernoulli-sets, ch:perfect-hash-filters.
% DEFINED: sec:cancellation, prop:complement-axis-swap, prop:mixed-intersection-escape.

\section{Cancellation and the Loss of One-Sidedness}
\label{sec:cancellation}

The operands of \Cref{sec:union-of-sets} and \Cref{sec:other-set-ops} were
drawn from one-sided sets on the perfect filter line. Some compositions kept
the result on the operands' axis; others did not. This section classifies the
three behaviors: a composition can stay on the operands' axis, walk to the
other axis, or leave the line entirely into the two-sided
interior\index{one-sided error}\index{two-sided error}. The
complement and the mixed-axis intersection drive the last two, and
\Cref{sec:perfect-filter-line} rereads the classification as a closure
question.

\subsection*{Complement swaps which axis}

\begin{proposition}[Complement swaps the perfect-filter axis]
\label{prop:complement-axis-swap}
If $\ASet{A}$ is one-sided, so is $\complement \ASet{A}$, on the other axis:
FNR-zero maps to FPR-zero and vice versa.
\end{proposition}

\begin{proof}
By \Cref{prop:complement-swap} the complement sends $(\fprate_A, \fnrate_A)$
to $(\fnrate_A, \fprate_A)$. If $\fnrate_A = 0 < \fprate_A$ then the image is
$(0, \fprate_A)$ with $\fprate_\complement = 0 < \fnrate_\complement$, which
is FPR-zero. The FPR-zero case is symmetric.
\end{proof}

A Bloom filter at $(0.01, 0)$ on the FNR-zero axis has complement $(0, 0.01)$
on the FPR-zero axis: the same information read dually, certain about absence
becoming certain about presence. The complement is a discrete reflection
across the rate plane, not a continuous deformation; it walks the line from
one arm to the other in a single step. Composing union with complement
illustrates this without leaving the line. Two FNR-zero filters at $\fprate_A
= 0.01$, $\fprate_B = 0.02$ have union $(0.0298, 0)$ by
\Cref{prop:union-fpr-fnr}, still FNR-zero, and complementing the union swaps
its rates to $(0, 0.0298)$ on the FPR-zero axis. The composed structure stays
one-sided; only the certain direction flips, from ``is $x$ in $A$ or $B$?''
answered with a sure $\bot$ to ``is $x$ outside both?'' answered with a sure
$\top$.

\subsection*{Mixed-axis intersection leaves the line}

The complement enters destructively the moment it sits inside an
intersection. Let $\ASet{A}$ be FNR-zero at $(0.01, 0)$ and $\ASet{B}$
FNR-zero at $(0.02, 0)$, and form the difference $\ASet{A} \setminus \ASet{B}
= \ASet{A} \cap \complement \ASet{B}$. The complement $\complement \ASet{B}$
sits at $(0, 0.02)$ on the FPR-zero axis, so the intersection pairs one
operand from each axis. By the corrected intersection formula
\eqref{eq:intersection-fpr}, the false-positive rate is \emph{not} the bare
product: with $\ASet{A} = (0.01, 0)$ and $\complement \ASet{B} = (0, 0.02)$,
the outside-both region survives,
\[
  \fprate_\cap \approx \fprate_A (1 - \fnrate_{\complement B})
              = 0.01 \times 0.98 = 0.0098,
  \qquad
  \fnrate_\cap = 1 - (1 - 0)(1 - 0.02) = 0.02.
\]
The composition lands at $\approx (0.0098, 0.02)$, in the two-sided interior,
not on either operand's axis. This is the same two-sidedness
\Cref{sec:other-set-ops} found for the difference directly, as it must be,
since the two are the same set\index{cancellation}.

\begin{proposition}[Mixed-axis intersection escapes the line]
\label{prop:mixed-intersection-escape}
Let $\ASet{A}$ be FNR-zero ($\fprate_A > 0$, $\fnrate_A = 0$) and $\ASet{B}$
FPR-zero ($\fprate_B = 0$, $\fnrate_B > 0$). Then $\ASet{A} \cap \ASet{B}$ is
two-sided: by \eqref{eq:intersection-fpr} its FPR is $w_2 \fprate_A(1 -
\fnrate_B) > 0$, contributed by the region outside both sets where $\ASet{A}$
false-positives and $\ASet{B}$ accepts, and by \eqref{eq:intersection-fnr}
its FNR is $\fnrate_B > 0$. The line is not preserved.
\end{proposition}

\begin{proof}
Substitute $(\fprate_A, 0)$ and $(0, \fnrate_B)$ into
\Cref{prop:intersection-fpr-fnr}. The FNR is $1 - (1 - 0)(1 - \fnrate_B) =
\fnrate_B$. In the FPR sum \eqref{eq:intersection-fpr} the first and third
terms carry a factor $\fprate_B = 0$ and vanish, leaving $w_2 \fprate_A(1 -
\fnrate_B)$, strictly positive whenever the outside-both region has positive
weight. Both output rates are positive, so the result is off both axes.
\end{proof}

The older intuition that this intersection ``lands on the FPR-zero axis'' is
an artifact of approximating $\fprate_\cap$ by the product $\fprate_A
\fprate_B = 0$; the linear term is what makes the escape immediate, with no
need for a leak in either operand's one-sidedness.

\subsection*{The three regimes}

The chapter-6 operations fall into three classes by their action on the
perfect filter line.

\emph{Axis-preserving.} Union and intersection of two FNR-zero sets stay
FNR-zero: $\fnrate = 0$ is preserved (\Cref{eq:union-fnr},
\Cref{eq:intersection-fnr} vanish), even though the surviving FPR is the
full three-region rate of §6.2. The FPR-zero side is symmetric.

\emph{Axis-swapping.} Complement of a one-sided set is one-sided on the other
axis (\Cref{prop:complement-axis-swap}): the result leaves each axis
individually but stays on the line as a whole.

\emph{Axis-leaving.} Mixed-axis intersection is two-sided
(\Cref{prop:mixed-intersection-escape}). Difference $\ASet{A} \setminus
\ASet{B} = \ASet{A} \cap \complement \ASet{B}$ is two-sided exactly when
$\ASet{A}$ and $\ASet{B}$ lie on the \emph{same} axis (so the complement
sends $\ASet{B}$ to the opposite axis and the intersection is mixed); it
stays one-sided only when the operands lie on opposite axes to begin with.
These are the compositions that exit the line into the interior, where both
directions of report are uncertain.

\subsection*{Why it matters}

Whether a composed set sits on the line governs the three later chapters that
take the line as their subject. \Cref{ch:classification-measures} simplifies
for a one-sided structure (an FNR-zero set has $\NPV = 1$ and one nonzero
rate to estimate) and exercises the full machinery for a two-sided one.
\Cref{ch:information-bounds} presents the classical floor $-\log_2 \fprate$
for FNR-zero sets and its two-sided form, and
\Cref{ch:optimal-bernoulli-sets} constructs a structure at the floor; a chain
that leaves the axis is priced by the two-sided form instead.
\Cref{ch:perfect-hash-filters} builds on the FPR-zero axis, and a mixed-axis
intersection with a Bloom filter destroys its certain $\top$. Composition
discipline, keeping a chain on one axis, is the price of staying near the
floor; \Cref{sec:perfect-filter-line} states the discipline as a closure
result.
